% CV de Estefanny Mejía
% Compila con pdfLaTeX (por ejemplo, en Overleaf).
\documentclass[10.5pt]{article}

\usepackage[T1]{fontenc}
\usepackage[utf8]{inputenc}
\usepackage[a4paper,margin=2cm,top=1.8cm]{geometry}
\usepackage[sc]{mathpazo}          % Palatino para títulos
\usepackage[scaled=0.92]{helvet}   % Helvetica para el texto
\renewcommand{\familydefault}{\sfdefault}
\usepackage{xcolor}
\usepackage{titlesec}
\usepackage{enumitem}
\usepackage{tabularx}
\usepackage[hidelinks]{hyperref}

% Colores del sitio web
\definecolor{tinta}{HTML}{4A2A33}
\definecolor{vino}{HTML}{8A4B5A}
\definecolor{oro}{HTML}{B97A6D}
\definecolor{suave}{HTML}{6E585D}
\definecolor{linea}{HTML}{ECC9CC}
\color{tinta}

\newcommand{\serif}[1]{{\fontfamily{ppl}\selectfont #1}}
\pagestyle{empty}
\setlength{\parindent}{0pt}
\raggedright
\setlist[itemize]{leftmargin=1.2em, itemsep=2pt, topsep=3pt, label={\color{oro}\textbullet}}

% Títulos de sección
\titleformat{\section}
  {\fontfamily{ppl}\selectfont\Large\color{vino}}
  {}{0pt}{}[{\color{linea}\titlerule[0.8pt]}]
\titlespacing*{\section}{0pt}{14pt}{7pt}

% Entrada: título, institución y fecha opcional a la derecha
\newcommand{\entrada}[3]{%
  \begin{tabularx}{\linewidth}{@{}X r@{}}
    \serif{\large #1} & {\small\color{oro}#3}\\
    {\color{suave}#2} & \\
  \end{tabularx}\vspace{4pt}}

\begin{document}

% Encabezado
\begin{center}
  {\fontfamily{ppl}\selectfont\fontsize{30}{34}\selectfont Estefanny Mejía}\\[4pt]
  {\small\color{oro}ECONOMISTA \,·\, DATA SCIENCE}\\[6pt]
  {\small\color{suave}\href{https://github.com/estefmej}{github.com/estefmej} \quad|\quad \href{https://estefmej.github.io}{estefmej.github.io}}
\end{center}

\section{Perfil}
Economista de la Universidad Continental con 6 años de experiencia. Trabajé como consultora en la
Organización Panamericana de la Salud (OPS). Actualmente curso el Diplomado de Data Science de la PUCP,
donde convierto datos públicos del Perú en análisis reproducibles, mapas y dashboards.

\section{Experiencia}
\entrada{Consultora}{Organización Panamericana de la Salud (OPS)}{}

\section{Formación}
\entrada{Diplomado en Data Science}{Pontificia Universidad Católica del Perú (PUCP)}{En curso}
\entrada{Economía}{Universidad Continental}{}

\section{Proyectos}
\serif{\large ¿Tener red pública garantiza agua segura en el Perú?}
\begin{itemize}
  \item Reconstruí el indicador ODS 6.1.1 con microdatos de la ENAPRES 2019--2025 y lo validé contra las cifras oficiales del MVCS y el INEI.
  \item El agua segura cayó de 41,5\,\% a 36,6\,\%; el cloro, no la red, es el cuello de botella. Dashboard interactivo por departamento.
\end{itemize}
\vspace{4pt}
\serif{\large El fuego y el agua}
\begin{itemize}
  \item Crucé focos de calor 2012--2026 con las zonas que producen agua y medí áreas quemadas con Sentinel-2.
  \item El 24\,\% de los focos de la temporada 2026 cayó en esas zonas; Puno concentra 4.011. Dashboard publicado en GitHub Pages.
\end{itemize}
\vspace{4pt}
\serif{\large Pobreza monetaria con la ENAHO 2024}
\begin{itemize}
  \item Pipeline automatizado de descarga, limpieza, gráficos y mapas que reproduce las cifras del INEI: 27,6\,\% nacional; 39,3\,\% rural y 24,8\,\% urbano.
\end{itemize}

\section{Herramientas}
Python \,·\, R \,·\, Power BI \,·\, Stata \,·\, QGIS \,·\, Streamlit \,·\, Mapas y GIS \,·\, Claude Code \,·\, Git y GitHub

\vspace{4pt}
{\color{suave}Datos con los que trabajo: microdatos del INEI (ENAHO, ENAPRES), imágenes satelitales (Sentinel-2) y focos de calor (NASA FIRMS, INPE).}

\end{document}
